\rhead{ML201: Fundamentals of AI \& Game Theory}
\phantomsection 
\section*{Fundamentals of AI \& Game Theory (ML201)}
\label{major:ML_1}

\noindent \small{\hyperref[main:scheme]{$\leftarrow$ Back to Scheme}} \vspace{0.5cm}

\noindent \textbf{Credits:} 3 (3-0-0)

\subsection*{Course Content}
\noindent\textbf{Unit 1} \\
\textit{Introduction to Artificial Intelligence:} History and foundations of AI, The Turing Test and Chinese Room Argument. Intelligent Agents: Agents and Environments, The Concept of Rationality, The Nature of Environments (Fully/Partially observable, Deterministic/Stochastic), Structure of Agents (Simple Reflex, Model-based, Goal-based, Utility-based).

\vspace{0.3cm}

\noindent\textbf{Unit 2} \\
\textit{Problem Solving by Search:} Problem-Solving Agents, Formulating Problems, Example Problems (8-Puzzle, TSP). Uninformed Search Strategies: Breadth-First Search (BFS), Depth-First Search (DFS), Depth-Limited Search, Iterative Deepening DFS. Informed (Heuristic) Search Strategies: Greedy Best-First Search, A* Search (Optimality and Completeness), Heuristic Functions (Admissibility and Consistency). Local Search Algorithms: Hill-climbing, Simulated Annealing.

\vspace{0.3cm}

\noindent\textbf{Unit 3} \\
\textit{Game Theory and Adversarial Search:} Games as Search Problems. Zero-Sum Games. Perfect Information Games: The Minimax Algorithm, Optimal decisions in games. Alpha-Beta Pruning (optimization). Imperfect Real-time decisions. Introduction to General Game Playing. Non-Zero-Sum Games: Prisoner's Dilemma, Nash Equilibrium, Cooperative vs. Non-Cooperative Games.

\vspace{0.3cm}

\noindent\textbf{Unit 4} \\
\textit{Knowledge Representation and Reasoning:} Logical Agents. Propositional Logic: Syntax and Semantics, Entailment, Inference rules. First-Order Logic (FOL): Syntax and Semantics, Quantifiers (Universal, Existential). Inference in FOL: Forward Chaining, Backward Chaining, and Resolution. Knowledge Engineering and Ontologies.

\vspace{0.3cm}

\noindent\textbf{Unit 5} \\
\textit{Uncertainty and Utility Theory:} Quantifying Uncertainty: Probability axioms, Random Variables, Independence, Bayes' Rule and its application. Making Simple Decisions: Utility Theory, Utility Functions, Decision Networks, The Value of Information. Introduction to Markov Decision Processes (MDP) basics (States, Actions, Rewards).

\newpage
\rhead{Curriculum Schema}